\chapter{\foreignlanguage{english}{DishBrain}}
\chaptermark{\foreignlanguage{english}{DishBrain}}
\label{sec:dishbrain}

\section{شبكة في طبق تلعب التنس}

\foreignlanguage{english}{DishBrain} (دماغ في طبق) هو الاسم الذي أطلقه الباحثون على منصة تلعب فيها مزرعة من عصبونات حية، مستنبتة على شريحة بأقطاب، لعبة تنس مبسّطة بمضرب واحد~\cite{Kagan2022}. وهي أكثر التجارب على الآلات المصنوعة من عصبونات حية إثارةً للنقاش (وعرض هذه الآلات في الفصل~\ref{sec:livingmachine})، وهي مناسبة لكتابنا خاصة. فالشبكة الصغيرة هنا متاحة كلها: كل دخل يحدده المجرِّب، وكل خرج يُسجَّل. وكل أسئلة الفصول السابقة، أي كيف تُرمَّز الإشارة، ومن يعلّم، وماذا يرى المسبار، يمكن أن تُطرح على شبكة لا عيون لها ولا عضلات ولا عصبونات \nt{DOPA}. فلنحلّل التجربة كما يحلّل المهندس منصة غيره: المخطط، والدخل، والخرج، والمعلّم، والنتيجة، والخلاف.

\section{المنصة}

تؤخذ العصبونات من قشرة جنين الفأر، نحو $800\,000$ خلية للشريحة، أو تُستنبت من خلايا جذعية بشرية، نحو مليون للشريحة. وتنمو على الشريحة بضعة أسابيع، فتتشابك زوائدها وتبدأ من تلقاء نفسها بإطلاق النبضات والرشقات. وتحتها، على مساحة $8\,\text{mm}^2$، $26\,000$ قطب من البلاتين؛ يمكن التسجيل من $1024$ منها في آن واحد، أما التحفيز ففي الممارسة عبر ثمانية.

وحول المزرعة حلقة مغلقة (الشكل~\ref{fig:dishloop}). يحاكي الحاسوب اللعبة: الكرة تطير وترتد عن الجدران، والمضرب يتحرك إلى أعلى وإلى أسفل. ويتحول موضع الكرة إلى تيار على ثمانية أقطاب ”حسية“. وتُعدّ النبضات من منطقتين ”حركيتين“ في نوافذ مدة كل منها $10\ms$، فتحرّك المضرب. وبعد كل إصابة أو إخفاق تتلقى المزرعة منبهًا آخر، هو التغذية الراجعة. ونوافذ $10\ms$ هي نبضة ساعة حاسوب المنصة لا المزرعة: فالشبكة نفسها تبقى تعمل عملًا لا متزامنًا، ككل شبكات هذا الكتاب.

\begin{figure}[t]
\centering
\begin{tikzpicture}[>=Stealth,
  blk/.style={draw,very thick,rounded corners=3pt,align=center,font=\small,minimum height=1.2cm},
  lab/.style={font=\scriptsize,align=center}]
\node[blk,fill=blue!6,text width=3.4cm] (C) at (0,0) {مزرعة على شريحة\\\scriptsize $\sim10^6$ عصبون، $26\,000$ قطب};
\node[blk,fill=orange!10,text width=3.0cm] (G) at (7.2,0) {الحاسوب:\\لعبة التنس};
\draw[->,very thick,blue!60!black] (G.north west) to[bend right=25] node[lab,above] {الكرة: \emph{المكان}: أيّ الأقطاب الثمانية،\\\emph{التردد}: $4\text{--}40$ نبضة/ث} (C.north east);
\draw[->,very thick,red!70!black] (C.south east) to[bend right=25] node[lab,below] {المضرب: عدّ نبضات المنطقتين\\”أعلى“ و”أسفل“ في $10\ms$} (G.south west);
\draw[->,thick,densely dashed,green!45!black] (G.east) -- ++(0.9,0) |- ++(0,-2.6) -| node[lab,pos=0.25,below] {إصابة: منبه متوقَّع؛ إخفاق: منبه غير متوقَّع} (C.south);
\end{tikzpicture}
\caption{حلقة \foreignlanguage{english}{DishBrain}. السهم الأزرق هو الدخل: موضع الكرة مرمَّز بترميز المكان وبترميز التردد (الفصل~\ref{sec:code}). والأحمر هو الخرج: الفرق بين عدّي المنطقتين يحرّك المضرب. والأخضر المتقطع هو التغذية الراجعة بعد كل جولة، وهي ”المعلّم“ الوحيد في المنصة.}
\label{fig:dishloop}
\end{figure}

\section{الدخل والخرج}

\emph{الدخل} مرمَّز بترميزين من الفصل~\ref{sec:code} معًا. أيّ الأقطاب الثمانية يُحفَّز يخبر على أي ارتفاع توجد الكرة: هذا ترميز المكان، كما في مستشعرات الفصل~\ref{sec:sensors}. وكم مرة تأتي نبضات التحفيز يخبر كم تبعد الكرة: $4$ في الثانية حين تكون عند الجدار البعيد، وأكثر خطيًا حتى $40$ في الثانية حين تكون عند جدار المضرب. وهذا ترميز التردد؛ وسعة نبضات الكرة $75\mV$. والمزرعة لا ”تعرف“ أيًّا من هذين الترميزين مسبقًا: الترميز حدده المجرِّب، وعلى الشبكة أن تتعلم قراءته.

\emph{الخرج} يُقرأ كما في واجهات الفصل~\ref{sec:debugging}. نبضات إحدى المنطقتين الحركيتين تدفع المضرب إلى أعلى، ونبضات الأخرى إلى أسفل؛ وفي أبسط صورة، في كل نافذة
\begin{empheq}[box=\keybox]{equation}
\Delta p \;=\; c\,\bigl(n_\uparrow - n_\downarrow\bigr),
\label{eq:dishreadout}
\end{empheq}
حيث $n_\uparrow$ و$n_\downarrow$ عددا نبضات المنطقتين في $10\ms$، و$c$ معامل المنصة. هذا هو الترميز بسلكين من الفصل~\ref{sec:glutcomputer}: اتجاه الحركة لا تحمله علامة الإشارة، بل أيّ السلكين أعلى صوتًا.

لننظر في ضجيج هذا الخرج. إذا أطلقت المنطقة ما مجموعه $R$ نبضة في الثانية، ففي النافذة $T=10\ms$ يكون متوسطها $RT$، وبحسب~\eqref{eq:rateerror} يكون التشتت النسبي $1/\sqrt{RT}$. عند $R=1000$ نبضة في الثانية يبلغ هذا نحو $30$ في المئة في كل نافذة؛ وعند $R=100$ أكثر من مئة. فالمضرب يرتجف حتمًا، ولا تستطيع الشبكة أن تتعلم إلا بطريقة واحدة: أن تزيح \emph{متوسط} الفرق بين العدّين في الاتجاه المطلوب. إنها القوانين نفسها التي تحدّ كل واجهة في الفصل~\ref{sec:debugging}، لكنها الآن على طرفي الحلقة كليهما.

\section{معلّم بلا دوبامين}

أطرف ما في المنصة معلّمها. فمزرعة عصبونات القشرة لا عصبونات \nt{DOPA} فيها، والمنصة لا ترسل أي إشارة من نوع ”أفضل مما توقعنا“. التغذية الراجعة من نوع آخر. إذا صدّت الشبكة الكرة، قُدّم إلى الأقطاب الحسية الثمانية كلها معًا منبه قصير \emph{متوقَّع}: نبضات بسعة $75\mV$، $100$ في الثانية مدة $0.1\,\text{s}$. وإذا أخفقت، استمر أربع ثوانٍ منبه يسميه الباحثون \emph{غير متوقَّع}: نبضات بسعة $150\mV$، $5$ في الثانية، في لحظات عشوائية وعلى أقطاب مختارة عشوائيًا؛ ثم أربع ثوانٍ من التوقف بلا تحفيز، وتنطلق الكرة من جديد~\cite{Kagan2022}.

انطلق الباحثون من فرضية أن الشبكة تتصرف بحيث يصير دخلها متوقَّعًا~\cite{Friston2010}. والمعنى عند المهندس بسيط: ليس لدى المزرعة إلا طريقة واحدة للتخلص من أربع ثوانٍ من الضجيج العشوائي، وهي أن تصدّ الكرة. وقد صادفنا الفكرة نفسها في الاقتصاد في النبضات في الفصل~\ref{sec:code} وفي التنبؤ بالإطار في الفصل~\ref{sec:recognition}.

لكن هل يلزم هنا التوقُّع بعينه؟ هناك آلية أخشن ممكنة أيضًا. لنفترض أن المنبه غير المتوقَّع بعد الإخفاق \emph{يخلط} فحسب التغيرات الأخيرة في الشبكة، وأن المنبه الهادئ بعد الإصابة لا يمسّها. لنصف ضبط الشبكة بمتجه واحد $\boldsymbol\psi$، ولتكن احتمالية إخفاق الشبكة في الضبط $\boldsymbol\psi$ هي $P_{\text{miss}}(\boldsymbol\psi)$. إذا كانت الدفعات متناظرة، أي أن الانتقال من $\boldsymbol\psi$ إلى $\boldsymbol\psi'$ محتمل احتمال الانتقال المعاكس، ففي النظام المستقر يتوازن تدفق المغادرة من كل ضبط مع تدفق القدوم إليه، ونسبة الوقت الذي تقضيه الشبكة في الضبط $\boldsymbol\psi$ هي
\begin{empheq}[box=\keybox]{equation}
\rho(\boldsymbol\psi) \;\propto\; \frac{1}{P_{\text{miss}}(\boldsymbol\psi)} .
\label{eq:dishdensity}
\end{empheq}
الضبط الذي يخفق أقل بعشر مرات يدوم أطول بعشر مرات. لا أحد يخبر الشبكة في أي اتجاه تغيّر أوزانها: الضبوط الجيدة تُهدم أقل فحسب.

اختبرنا ذلك على نموذج يختزل المزرعة في عددين: يصل المضرب إلى النقطة $a\,y+b$ حين تصل الكرة إلى الارتفاع $y$، ويصدّها إذا كان خطؤه أقل من $0.1$. بعد الإخفاق يتلقى العددان دفعة عشوائية، وبعد الإصابة يبقيان على حالهما. فارتفعت نسبة الكرات المصدودة في $2000$ جولة من $0.21$ في أول مئتين إلى $0.38$ في آخر خمسمئة (أربعون ”مزرعة“، الشكل~\ref{fig:dishmodel}). وإذا أُعطيت الدفعة بعد \emph{كل} جولة، فلا معلومة في التغذية الراجعة، وتبقى النسبة نحو $0.12$؛ وإذا لم تكن دفعات أصلًا ثبتت عند $0.19$. والتجربة على المنصة تتفق مع هذا: لم يظهر نمو ذو دلالة داخل الجلسة إلا مع التغذية الراجعة الكاملة. ومع الصمت بدل المنبهات لعبت المزرعة في آخر الجلسة مع ذلك أفضل مما بلا تغذية راجعة، لكن بأثر أصغر~\cite{Kagan2022}؛ ونلاحظ أن في هذا الشرط أيضًا يأتي بعد الإخفاق توقف طويل وبعد الإصابة توقف قصير، فلا يختفي الفرق بين النتيجتين كليًا.

\begin{figure}[t]
\centering
\begin{tikzpicture}[x=1cm,y=5cm]
\draw[->,gray!70!black] (0,0) -- (10.6,0);
\draw[->,gray!70!black] (0,0) -- (0,0.5) node[above,font=\small,black] {نسبة المصدود};
\foreach \v in {0.1,0.2,0.3,0.4}{\draw[gray!60!black] (-0.06,\v) -- (0.06,\v); \node[left,font=\scriptsize] at (0,\v) {\v};}
\foreach \x/\f/\l/\lab in {1.8/0.21/0.38/{ضجيج بعد\\الإخفاق},5.2/0.13/0.11/{ضجيج بعد\\كل جولة},8.6/0.19/0.19/{بلا ضجيج}}{
  \fill[blue!35] (\x-0.8,0) rectangle (\x-0.05,\f);
  \fill[red!60!black] (\x+0.05,0) rectangle (\x+0.8,\l);
  \node[below,font=\scriptsize,align=center] at (\x,-0.01) {\lab};
  \node[above,font=\scriptsize] at (\x-0.425,\f) {\f};
  \node[above,font=\scriptsize] at (\x+0.425,\l) {\l};}
\fill[blue!35] (7.4,0.47) rectangle (7.7,0.495); \node[right,font=\scriptsize] at (7.75,0.482) {أول 200};
\fill[red!60!black] (7.4,0.41) rectangle (7.7,0.435); \node[right,font=\scriptsize] at (7.75,0.422) {آخر 500};
\end{tikzpicture}
\caption{نموذج ”الخلط عند الإخفاق“ (\foreignlanguage{english}{\texttt{dishbrain\_model.py}}): نسبة الكرات المصدودة في بداية $2000$ جولة وفي نهايتها، بالمتوسط على أربعين مزرعة نموذجية. لا يحدث التعلم إلا إذا جاءت الدفعة بعد الإخفاق لا بعد الإصابة، كما في التجربة، حيث لم يظهر نمو ذو دلالة داخل الجلسة إلا مع التغذية الراجعة الكاملة.}
\label{fig:dishmodel}
\end{figure}

\begin{keyblock}
\begin{proposition}[المعلّم الخالط]
\label{prop:dishscramble}
إذا كان الفشل يهزّ الشبكة والنجاح يتركها وشأنها، انجرفت الشبكة من تلقاء نفسها نحو الضبوط التي تخطئ أقل: الوقت الذي تقضيه في ضبط ما يتناسب عكسيًا مع احتمالية الفشل~\eqref{eq:dishdensity}. ولا يحتاج هذا التعلم إلى إشارة مكافأة ولا إلى علامتها ولا إلى التوقُّع بحد ذاته؛ يكفي أن تختلف التغذية الراجعة بعد النجاح عنها بعد الفشل. تغيّر سلوك المزرعة مُقاس؛ أما الآلية العاملة في الطبق، أهي التنبؤ بالدخل أم الخلط، فلم تُحدَّد.
\end{proposition}
\end{keyblock}

قارن بالفصل~\ref{sec:dofa}. هناك أيضًا كان الضجيج أداة بحث، لكن الدوبامين كانت له علامة: الخطأ $\delta$ كان يخبر في أي اتجاه تُزاح الأوزان، وكانت الخطوة تسير على التدرج~\eqref{eq:perturbation}. أما هنا فلا علامة: ليس إلا ”أبقِ“ أو ”اهزز“. وهذا البحث أخشن وأبطأ، لكنه يطلب من الشبكة أقل: لا عصبونات \nt{DOPA}، ولا ناقد، ولا أثر أهلية ممتد على ثوانٍ.

\section{ما قيس وما الخلاف عليه}

دامت كل لعبة $20$ دقيقة؛ وقورنت أول $5$ دقائق بآخر $15$. في المزارع ذات التغذية الراجعة الكاملة طالت سلاسل الكرات المصدودة وقلّت الإخفاقات الفورية؛ ولم يُرَ شيء من هذا في التجارب الضابطة بلا خلايا، وبلا لعبة، وبـ”مضرب“ يحرّكه الحاسوب. ومزارع الخلايا البشرية صدّت في المتوسط مددًا أطول من مزارع الفأر. والتحسن متين إحصائيًا، على تسع مزارع من الفأر وإحدى عشرة بشرية، في أكثر من مئة لعبة لكل مجموعة، لكنه صغير: لم تصر الشبكة لاعبًا جيدًا، بل صارت أفضل قليلًا من العشوائية. وهو متين داخل الجلسة؛ أما التعلم الثابت من جلسة إلى أخرى على مدى أيام فلم يحصل عليه الباحثون~\cite{Kagan2022}. ووجدت الدراسة التالية للمجموعة نفسها أن فعالية المزرعة في أثناء اللعب تقترب من النظام \emph{الحرج}، أي الحد بين الخمود والانهيار المتسلسل، وهو العيش على الحافة نفسه الذي في الفصل~\ref{sec:gaba}، وكلما اقتربت كان اللعب أفضل~\cite{Habibollahi2023}.

الخلاف الرئيسي لم تُثِره الأرقام، بل الكلمات. فعنوان المقالة زعم أن العصبونات في الطبق ”تُظهر وعيًا“ (\foreignlanguage{english}{sentience})، واعترضت مجموعة من الباحثين: تغيّر السلوك في حلقة مغلقة ليس ذكاءً ولا إحساسًا بعد، والاستنتاجات يجب أن تُصاغ بتواضع أكبر~\cite{Balci2023}. ومن الزاوية الهندسية سؤالان مفتوحان، وكلاهما عن الآلية. الأول: هل تتعلم الشبكة، أي هل تتغير أوزانها تغيرًا دائمًا، أم أن التغذية الراجعة تزيح حالتها الراهنة فحسب؟ والثاني: هل يلزم التوقُّع بعينه، أم يكفي أي فرق بين الاستجابة للنجاح والاستجابة للفشل، كما في القضية~\ref{prop:dishscramble}؟ والمنصة التي يُتاح فيها كل قطب تسمح بالإجابة عن السؤالين معًا، ولعل هذه أهم قيمها.

\section{مواصفات المنصة: كيف نعيد بناءها}
\label{sec:dishspec}

جمعنا أدناه كل ما يلزم لبناء منصة كهذه من جديد: العتاد، والحلقة، وترميز الدخل، وقراءة الخرج، والتغذية الراجعة، وبروتوكول التجربة. وكل معامل موسوم بمصدره. ”المقالة“: القيمة مأخوذة من نص الدراسة~\cite{Kagan2022}. ”اختيار“: القيمة غير موجودة في المقالة، وعلى من يعيد التجربة أن يحددها بنفسه؛ ونقترح قيمة افتراضية ونبيّن كيف تُختبر.

\subsection*{العتاد والمزرعة}

\begin{center}
\footnotesize
\setlength{\tabcolsep}{4pt}
\renewcommand{\arraystretch}{1.15}
\begin{tabular}{>{\raggedright}p{4.0cm}>{\raggedright}p{6.9cm}>{\raggedright\arraybackslash}p{1.6cm}}
\toprule
المعامل & القيمة & المصدر \\
\midrule
الشريحة & مصفوفة \foreignlanguage{english}{MaxOne}: $26\,000$ قطب من البلاتين على مساحة $8\,\text{mm}^2$ & المقالة \\
التسجيل & حتى $1024$ قناة في آن واحد، $20\,000$ عينة في الثانية & المقالة \\
التحفيز & حتى $32$ قطبًا تقنيًا، وفي التجربة $8$ مستقلة & المقالة \\
الخلايا & قشرة جنين الفأر؛ عصبونات بشرية من خلايا جذعية (طريقتا استنبات)؛ خلايا \foreignlanguage{english}{HEK293T} (ليست عصبونات) للمقارنة الضابطة & المقالة \\
الزرع & نحو $10^6$ خلية للشريحة & المقالة \\
عمر المزرعة عند التجربة & الفأر: من $\sim10$ أيام؛ الإنسان: $14\text{--}24$ يومًا أو $\sim80$ يومًا بحسب طريقة الاستنبات & المقالة \\
\bottomrule
\end{tabular}
\end{center}

\subsection*{الحلقة والزمن}

تعمل الحلقة بنوافذ مدة كل منها $10\ms$ ($200$ عينة): في كل نافذة تُعدّ النبضات على أقطاب المنطقتين الحركيتين، وتُحدَّث اللعبة، ويُصدر التحفيز التالي. والتأخير من النبضة في المزرعة إلى المنبه المجيب نحو $5\ms$. وفي أثناء اللعب تمرّ إشارة كل قطب بمرشح بِسل تمرير عالٍ من الرتبة الثانية بتردد قطع $100\,\text{Hz}$؛ وتُنعَّم القيمة المطلقة للإشارة بمرشح بِسل تمرير منخفض من الرتبة الأولى بتردد $1\,\text{Hz}$، وعتبة النبضة متناسبة مع هذه القيمة المنعَّمة. أما عتبة ستة انحرافات معيارية للخلفية فتذكرها المقالة لقياسات منفصلة للفعالية، لا للعب. وكي لا يُحسب التحفيز استجابةً من المزرعة، تُخمد كل الإشارات حين يصل في آن واحد أكثر من $15$ نبضة كبيرة، تزيد على $75\mV$. كل هذا من المقالة.

\subsection*{ترميز الكرة}

\begin{center}
\footnotesize
\setlength{\tabcolsep}{4pt}
\renewcommand{\arraystretch}{1.15}
\begin{tabular}{>{\raggedright}p{3.4cm}>{\raggedright}p{7.5cm}>{\raggedright\arraybackslash}p{1.6cm}}
\toprule
المعامل & القيمة & المصدر \\
\midrule
المنطقة الحسية & $8$ أقطاب مرتبة بحيث تعني الأقطاب المتجاورة مواضع متجاورة للكرة & المقالة \\
ترميز المكان & رقم القطب $k=1,\dots,8$ يحدد موضع الكرة بالنسبة إلى مركز المضرب & المقالة \\
أي إحداثي & الموضع الرأسي للكرة؛ وللإعادة: بالنسبة إلى المضرب، $k=\lceil 8(y-p+\tfrac12)\rceil$ عند $y-p\in[-\tfrac12,\tfrac12]$ & المقالة؛ والصيغة اختيار \\
ترميز التردد & تردد التحفيز من $4$ إلى $40$ نبضة/ث يحمل المسافة إلى المضرب & المقالة \\
اتجاه السُّلَّم & $4$ نبضات/ث حين تكون الكرة عند الجدار المقابل، وخطيًا حتى $40$ نبضة/ث حين تكون عند جدار المضرب: $f=4+36\,(1-x/L)$، حيث $x$ المسافة إلى جدار المضرب و$L$ عرض الملعب & المقالة \\
شكل النبضة & مستطيلة قصيرة ثنائية الطور: الطور الموجب أولًا، ثم السالب & المقالة \\
السعة & $75\mV$؛ اختيرت كي لا تُجبر الخلايا المثبَّطة على الإطلاق & المقالة \\
مدة الطورين & غير مذكورة؛ تؤخذ الإعدادات القياسية لنظام التحفيز ولا تُغيَّر في أثناء التجربة & اختيار \\
\bottomrule
\end{tabular}
\end{center}

\subsection*{قراءة المضرب}

في المقالة منطقتان حركيتان: النبضات في الأولى تحرّك المضرب إلى أعلى، وفي الثانية إلى أسفل؛ وتُوزن مساهمة المنطقتين لمراعاة اختلاف كثافة الخلايا واختلاف الفعالية~\cite{Kagan2022}. والصيغة الدقيقة غير مذكورة. وللإعادة نأخذ القاعدة~\eqref{eq:dishreadout}، $\Delta p=c\,(n_\uparrow-n_\downarrow)$ لكل نافذة $10\ms$، بوزن يسوّي بين المنطقتين بحسب متوسط فعاليتهما في الراحة (اختيار): يُقسم $n_\uparrow$ و$n_\downarrow$ على متوسط عدّ منطقتيهما في النافذة، مقيسًا قبل اللعب. ونختار المعامل $c$ بحيث يزيح فرقٌ بمقدار عدّ متوسط واحد المضربَ بنسبة $1\%$ من ارتفاع الملعب في النافذة (اختيار)؛ وعندئذ يقطع المضرب الملعب كله في $1\,\text{s}$ إذا رجحت إحدى المنطقتين رجحانًا ثابتًا.

ومواقع المنطقتين على الشريحة لم تُحدَّد في نص المقالة بإحداثيات؛ ليس هناك إلا مخطط. وللإعادة تكفي ثلاث مجموعات غير متداخلة من الأقطاب: ثمانية أقطاب حسية في الوسط، ومجموعة أقطاب تسجيل على كل جانب، إحداهما ”أعلى“ والأخرى ”أسفل“ (اختيار).

\subsection*{اللعبة}

أبعاد الملعب وطول المضرب وسرعة الكرة غير مذكورة في المقالة؛ لا يُذكر إلا أن الكرة تطير بسرعة ثابتة وترتد عن الجدران. وللإعادة (اختيار): ملعب $1\times1$، ومضرب على الجدار الأيمن طوله $0.2$ (إصابة عند $|y-p|<0.1$، كما في النموذج~\foreignlanguage{english}{\texttt{models/dishbrain\_model.py}})، والكرة تقطع الملعب في $1\,\text{s}$. والإخفاق من غير أي كرة مصدودة في الجولة يُعدّ ”إرسالًا ساحقًا“ (\foreignlanguage{english}{ace})؛ والجولة التي تزيد على ثلاث إصابات متتالية تُعدّ ”طويلة“ (المقالة).

\subsection*{التغذية الراجعة}

\begin{center}
\footnotesize
\setlength{\tabcolsep}{4pt}
\renewcommand{\arraystretch}{1.15}
\begin{tabular}{>{\raggedright}p{2.6cm}>{\raggedright}p{8.3cm}>{\raggedright\arraybackslash}p{1.6cm}}
\toprule
الحدث & ما يُقدَّم & المصدر \\
\midrule
إصابة & منبه متوقَّع: الأقطاب الحسية الثمانية كلها في آن واحد، $75\mV$، $100$ نبضة/ث، $100\ms$؛ ويُقطع تحفيز الكرة المعتاد خلال ذلك & المقالة \\
إخفاق & منبه غير متوقَّع: $150\mV$، $5$ نبضات/ث، $4\,\text{s}$، في لحظات عشوائية وعلى أقطاب مختارة عشوائيًا من الثمانية؛ ثم $4\,\text{s}$ بلا تحفيز، وتنطلق الكرة في اتجاه عشوائي & المقالة \\
قانون العشوائية & غير مذكور؛ وللإعادة: لحظات النبضات تدفق بواسوني بمتوسط تردد $5$ نبضات/ث & اختيار \\
\bottomrule
\end{tabular}
\end{center}

\subsection*{البروتوكول والتجارب الضابطة}

تدوم الجلسة الواحدة $20$ دقيقة؛ وتُقارن أول $5$ دقائق بآخر $15$ (المقالة). ومقاييس النتيجة ثلاثة: متوسط طول الجولة، ونسبة الإرسالات الساحقة، ونسبة الجولات الطويلة. وشروط التجربة (المقالة):
\begin{itemize}
\item \emph{المنبه}: الحلقة الكاملة، كما وُصفت أعلاه؛
\item \emph{الصمت}: بدل منبه الإصابة أو الإخفاق المدة نفسها بلا أي تحفيز، ثم تنطلق الكرة في اتجاه عشوائي؛
\item \emph{بلا تغذية راجعة}: بعد الإخفاق ترتد الكرة فحسب وتواصل طيرانها، ويستمر تحفيز موضع الكرة؛
\item \emph{الراحة}: اللعبة تجري والمضرب يتحرك بحسب الفعالية، لكن لا تحفيز أصلًا؛
\item \emph{الضوابط}: شريحة فيها وسط الاستنبات وحده؛ خلايا \foreignlanguage{english}{HEK293T}؛ ”مزرعة في الحاسوب“، أي محاكاة بدل الخلايا.
\end{itemize}
حجم العينة في السلسلة الرئيسية: $9$ مزارع من الفأر ($101$ جلسة)، و$11$ مزرعة بشرية ($138$)، و$6$ شرائح بوسط الاستنبات ($80$)، و$20$ مزرعة في الراحة ($42$)، و$3$ محاكيات ($38$)، والمجموع $399$ جلسة. وفي سلسلة التغذية الراجعة $486$ جلسة على مدى $3$ أيام بمعدل $3$ جلسات يوميًا، بالشروط ”المنبه“ و”الصمت“ و”بلا تغذية راجعة“ ($n=27$ و$17$ و$15$). ولم تظهر زيادة متوسط طول الجولة من أول $5$ دقائق إلى آخر $15$ إلا في المزارع الحية. وفي سلسلة التغذية الراجعة لم يظهر نمو ذو دلالة مع الزمن إلا في شرط ”المنبه“؛ وفي آخر الجلسة تفوّق شرط ”الصمت“ مع ذلك على ”بلا تغذية راجعة“، لكن بأثر أصغر، ولم يكن هناك تعلّم ثابت من جلسة إلى أخرى على مدى ثلاثة أيام~\cite{Kagan2022}.

\subsection*{دورة المنصة خطوةً خطوة}

كل $10\ms$ يفعل حاسوب المنصة ما يلي.
\begin{enumerate}
\item يقرأ عدد النبضات $n_\uparrow$ و$n_\downarrow$ في المنطقتين الحركيتين خلال النافذة المنقضية، ويقسمهما على متوسط عدّ المنطقة في الراحة.
\item يزيح المضرب بمقدار $\Delta p=c\,(n_\uparrow-n_\downarrow)$ ويحصره بين حدّي الملعب.
\item يحرّك الكرة خطوة واحدة؛ ويعكسها عن الجدران العلوي والسفلي والأيسر.
\item إذا بلغت الكرة الجدار الأيمن: عند $|y-p|<0.1$ فهي إصابة، ويزيد عدّاد الجولة واحدًا، ويُقدَّم منبه الإصابة ($100\ms$)؛ وإلا فهي إخفاق، وتُسجَّل الجولة، ويُقدَّم منبه الإخفاق ($4\,\text{s}$)، ثم توقف $4\,\text{s}$، وتنطلق الكرة من جديد.
\item وإلا يختار القطب $k$ بحسب الموضع الرأسي للكرة والتردد $f$ بحسب المسافة إلى جدار المضرب، ويقدّم إلى القطب $k$ نبضات ثنائية الطور بسعة $75\mV$ بالتردد $f$ حتى النافذة التالية.
\end{enumerate}
وهذا يكفي لبناء منصة متوافقة مع المنشورة، ولاختبار السؤال الرئيسي في الفصل: هل يعلّم المزرعةَ التوقُّعُ، أم مجرد الفرق بين الاستجابة للإصابة والاستجابة للإخفاق؟ ولهذا يحسن أن يُضاف إلى شروط المقالة الثلاثة شرط رابع ليس فيها: بعد الإخفاق منبه \emph{متوقَّع} بالمدة والطاقة نفسيهما اللتين للمنبه غير المتوقَّع. فإن تعلمت المزرعة مع ذلك، فالأمر في الفرق لا في التوقُّع.


\section{خلاصة}

\begin{table}[t]
\centering
\footnotesize
\setlength{\tabcolsep}{4pt}
\renewcommand{\arraystretch}{1.15}
\begin{tabular}{>{\raggedright}p{2.1cm}>{\raggedright}p{4.2cm}>{\raggedright}p{1.9cm}>{\raggedright\arraybackslash}p{4.6cm}}
\toprule
عقدة المنصة & كيف بُنيت & فصل الكتاب & ما هو معروف \\
\midrule
الدخل & المكان (أيّ الأقطاب الثمانية) والتردد ($4\text{--}40$ نبضة/ث) & \ref{sec:code}، \ref{sec:sensors} & يحدده المجرِّب \\
الخرج & الفرق بين عدّي المنطقتين في $10\ms$~\eqref{eq:dishreadout} & \ref{sec:glutcomputer}، \ref{sec:debugging} & يحدده المجرِّب؛ والضجيج وفق~\eqref{eq:rateerror} \\
المعلّم & منبه متوقَّع بعد الإصابة، وغير متوقَّع بعد الإخفاق؛ ولا دوبامين & \ref{sec:dofa} & نمو ذو دلالة داخل الجلسة مع التغذية الراجعة الكاملة فقط؛ ومع الصمت أثر أصغر؛ وغير ثابت من جلسة إلى أخرى \\
الآلية & التنبؤ بالدخل أو الخلط عند الإخفاق~\eqref{eq:dishdensity} & \ref{sec:dofa}، \ref{sec:recognition} & لم تُحدَّد؛ وكلتاهما تفسّر النتيجة \\
نظام الشبكة & الاقتراب من النظام الحرج في أثناء اللعب & \ref{sec:gaba} & صلته بجودة اللعب مقيسة \\
”الوعي“ & قول الباحثين & --- & لم يُبيَّن؛ ومطعون فيه \\
\bottomrule
\end{tabular}
\caption{\foreignlanguage{english}{DishBrain} بعين المهندس.}
\label{tab:dishbrain}
\end{table}

\foreignlanguage{english}{DishBrain} آلة من عصبونات حية في أبسط صورها: الدخل مرمَّز بترميز، والخرج مقروء بالعدّ، والمعلّم فرق بين الاستجابة للنجاح والاستجابة للفشل (الجدول~\ref{tab:dishbrain}). شبكة بلا عيون ولا عضلات ولا دوبامين تعلّمت اللعب قليلًا جدًا، وهذا وحده يجعلها منصة اختبار لأفكار الكتاب. وأيًّا كانت الآلية، تنبؤًا أو خلطًا أو شيئًا ثالثًا، فسيوجد الجواب كما ينصح الفصل~\ref{sec:debugging}: بالمسبار، وبإشارة الاختبار، وبتعطيل الأجزاء، على آلة نراها أخيرًا كاملة.

\section{تمارين}

\begin{exercise}[\lvlA]
\label{ex:22:1}
\emph{كم يلزم أن يُزاح العدّ.} تطلق المنطقتان معًا $R_\uparrow+R_\downarrow=2000$ نبضة في الثانية، والتدفقات بواسونية. (أ)~ما التشتت النسبي لعدّ منطقة واحدة في $10\ms$ عند $R=1000$ و$R=100$؟ (ب)~في $1\,\text{s}$ من طيران الكرة تتجمع الإزاحات~\eqref{eq:dishreadout}. ما أصغر $R_\uparrow-R_\downarrow$ يكون عندها متوسط الإزاحة ضعف التشتت؟
\end{exercise}

\begin{exercise}[\lvlA]
\label{ex:22:2}
\emph{كم جولة يلزم لرؤية التعلم.} الجولات مستقلة، ونسبة المصدود ثنائية الحدين. كم جولة يلزم في كل من الفترتين كي تتجاوز زيادة نسبة المصدود انحرافين معياريين للفرق: (أ)~من $0.20$ إلى $0.25$؛ (ب)~من $0.21$ إلى $0.38$، كما في النموذج؟
\end{exercise}

\begin{exercise}[\lvlB]
\label{ex:22:3}
\emph{اشتقاق~\eqref{eq:dishdensity}.} الضبط $\boldsymbol\psi$ يجري على مجموعة منتهية؛ عند الإخفاق (باحتمالية $P(\boldsymbol\psi)>0$) تنتقل الشبكة إلى $\boldsymbol\psi'$ باحتمالية متناظرة $K(\boldsymbol\psi,\boldsymbol\psi')$، وعند الإصابة تبقى. (أ)~برهن أن $\rho\propto1/P$ توزيع مستقر. (ب)~ما نسبة الإخفاقات مع ضبطين فيهما $P=0.9$ و$P=0.1$؟
\end{exercise}

\begin{exercise}[\lvlB]
\label{ex:22:4}
\emph{المنطقة الماصّة في النموذج.} يصل المضرب إلى $p=\min\bigl(1,\max(0,ay+b)\bigr)$، والكرة إلى الارتفاع $y\sim U(0,1)$، والإصابة إذا $|p-y|<h=0.1$. (أ)~برهن أن الشبكة تصدّ كل الكرات عند $|b|<h$ و$|a-1+b|<h$، وأوجد مساحة هذه المنطقة. (ب)~أوجد عدديًا متوسط نسبة المصدود عند $a\sim U(-1,1)$ و$b\sim U(0,1)$.
\end{exercise}

\begin{exercise}[\lvlB]
\label{ex:22:5}
\emph{محاكاة: سياج حول الضبوط.} في نسخة من \foreignlanguage{english}{\texttt{dishbrain\_model.py}} أعطِ $200$ مزرعة $20\,000$ جولة لكل منها. ما نسبة المصدود في آخر $500$؟ وماذا لو قُصّ الضبط بعد كل دفعة إلى المستطيل $a\in[-1,2]$، $b\in[-1,1]$؟ فسّر الفرق بالتمرين~\ref{ex:22:3}.
\end{exercise}

\begin{exercise}[\lvlB]
\label{ex:22:6}
\emph{تصميم ترميز الدخل.} تُصدّ الكرة إذا كان خطأ الارتفاع أقل من $h=0.1$؛ تقرأ الشبكة الدخل في $T=0.25\,\text{s}$، وحتى التردد $r$ يمكن تمييز نحو $\sqrt{rT}$ مستوى، والتحفيز لا يتجاوز $40$ نبضة في الثانية. (أ)~هل يكفي ترميز المكان على ثمانية أقطاب؟ (ب)~هل يمكن نقل الارتفاع بالتردد على قطب واحد؟
\end{exercise}

\begin{exercise}[\lvlC]
\label{ex:22:7}
\emph{بحث: تنبؤ أم خلط؟} عمّم النموذج على $d$ عددًا قابلًا للضبط ($p=\sum_{i<d}a_iy^i$). كيف يزداد مع $d$ عدد الجولات حتى نسبة مصدود $0.5$ عند الخلط وعند القاعدة ذات علامة الخطأ~\eqref{eq:perturbation}؟ وأيّ تجربة على المنصة تفصل بين الفرضيتين؟
\end{exercise}
